% !TEX program = xelatex
% configure agent enhancement analysis and OMX comparison report.
\documentclass[UTF8,a4paper,11pt]{ctexart}

\usepackage[a4paper,margin=2.05cm,headheight=14pt]{geometry}
\usepackage{amsmath,amssymb,mathtools}
\usepackage{booktabs,longtable,array,tabularx,multirow}
\usepackage{graphicx}
\usepackage{xcolor}
\usepackage{listings}
\usepackage{tikz}
\usetikzlibrary{arrows.meta,positioning,fit,calc}
\usepackage{url}
\usepackage{hyperref}
\usepackage{microtype}

\definecolor{primary}{RGB}{23,54,93}
\definecolor{accent}{RGB}{31,119,180}
\definecolor{evidence}{RGB}{30,132,73}
\definecolor{warning}{RGB}{197,90,17}
\definecolor{muted}{RGB}{91,101,115}
\definecolor{panel}{RGB}{244,247,250}
\hypersetup{colorlinks=true,linkcolor=primary,urlcolor=accent,pdftitle={configure Agent Enhancement 与 OMX 对照分析}}
\urlstyle{tt}
\newcommand{\code}[1]{\path{#1}}
\newcommand{\pathref}[1]{\path{#1}}
\newcommand{\verified}{\textcolor{evidence}{\textbf{确证}}}
\newcommand{\inferred}{\textcolor{accent}{\textbf{推断}}}
\newcommand{\unverified}{\textcolor{warning}{\textbf{未验证}}}
\newcommand{\metric}[1]{\textcolor{primary}{\textbf{#1}}}
\setlength{\parindent}{2em}
\setlength{\parskip}{0.35em}
\setlength{\emergencystretch}{3em}
\lstset{basicstyle=\ttfamily\small,backgroundcolor=\color{panel},frame=single,
  breaklines=true,breakatwhitespace=false,keepspaces=true,columns=flexible,
  keywordstyle=\color{primary}\bfseries,commentstyle=\color{evidence!70!black},
  stringstyle=\color{accent},numbers=left,numberstyle=\tiny\color{muted}}
\tikzset{
  box/.style={draw=primary,fill=primary!6,rounded corners=2pt,align=center,
    text width=3.05cm,minimum height=0.9cm,font=\small},
  store/.style={draw=evidence,fill=evidence!7,rounded corners=2pt,align=center,
    text width=3.25cm,minimum height=0.9cm,font=\small},
  arrow/.style={-{Latex[length=2mm]},draw=primary,semithick},
  relation/.style={-{Latex[length=1.7mm]},draw=muted,dashed,thin}
}

\title{configure Agent Enhancement 与 OMX 的\\结构分析、优化与对照测试}
\author{configure 工程分析记录}
\date{2026 年 9 月 9 日}

\begin{document}
\maketitle

\begin{abstract}
本文分析 \pathref{/Users/zhangxin/configure} 中面向 Codex、OpenCode 和 Claude
的 agent 增强部分，并与同一工作目录中引入的 OMX 项目
\pathref{refer/git-rep/oh-my-codex} 对照。工作对象的优化集中在统一启动器、跨平台运行时、会话恢复、分层上下文发现、生命周期事件、停止边界、并发锁、hook 适配和只读状态查询。结论按证据等级书写：configure 在低依赖、三类 CLI 统一入口、Unix/Windows 对齐、运行上限和可观测性方面形成了可复现优势；OMX 在多 worker 编排、tmux pane 权威、动态 scaling、worktree、角色路由、原子状态和大规模 hook 扩展方面仍明显更完整。因而本文不把轻量部署优势外推为整体平台超越，并给出已验证的优势范围与剩余差距。
\end{abstract}

\tableofcontents

\section{目的与前期准备}
\label{sec:purpose}
本次任务的对象是 configure 项目中的 agent enhancement，目标包括理解对照对象的优势、定位工作对象的劣势、实现针对性优化、运行对照测试，并生成可追溯报告。报告参考
\pathref{/Users/zhangxin/PyQCU/docs/report_multigrid_quda_pyqcu_20260906.tex} 的 A4 学术报告结构，采用 XeLaTeX 编译。分析主体属于代码/软件工作，因此使用 15 步代码分析框架；物理模型不属于本任务范围。

前置环境为 macOS、Bash、Git、Node/npm 和 XeLaTeX。configure 没有构建框架，Shell 校验以
\code{bash -n} 为主，agent 行为用 fake CLI 回归脚本覆盖。OMX 的 TypeScript 构建需要 Node 20 及以上；近期回归套件还需要 Rust/Cargo。\verified\ 这些条件分别由 configure 根部
\pathref{AGENTS.md:36--45}、OMX 的 \pathref{package.json:1--70} 和实际命令结果支持。

\section{输入是什么}
输入分为四组。第一组是工作代码：\code{bin/agent.sh}、\code{bin/agent.bat}、共享运行时、状态查询脚本、hook 适配层及回归测试。第二组是运行参数：launcher 名称、模型、reasoning effort、OpenCode variant、输入文件、继续间隔、回合上限、运行时长上限、停止文件和 resume run id。第三组是持久状态：\code{data/runs/<run-id>/} 下的 manifest、state、事件、上下文、日志和输入清单。第四组是对照证据：OMX HEAD
\code{0a43c09b1c4d7dbd03cfa7d9c4c42738d5d31584}、版本 \code{0.21.4}、源码与测试命令。

运行时输入的边界由 \code{bin/agent.sh:250--280}、\code{bin/agent.sh:535--570} 和
\code{bin/agent.sh:912--952} 的参数解析共同定义；持久化字段由
\code{bin/agent-runtime.sh:62--104} 定义。\inferred\ 将所有 launcher 的输入收敛到同一个运行记录，是为了让外部查询不再依赖某一种 CLI 的原始日志格式。

\section{输出应该是什么}
预期输出有三层：

\begin{enumerate}
  \item 用户层输出：TUI 或 headless agent 调用、紧凑的实时日志、错误状态和退出码。
  \item 机器层输出：\code{manifest.env}、\code{state.env}、\code{events.jsonl}、\code{context.txt}、\code{agent.log} 和 \code{inputs.txt}；这些文件不执行、不进入工作树提交。
  \item 交付层输出：测试结果、比较结论、\code{docs/analy_agent_configure_20260909.tex} 和同名 PDF，以及符合 tag 约定的 Git 版本记录。
\end{enumerate}

\code{data/README.md:7--27} 明确了数据文件语义、默认路径和 Git 边界；
\code{bin/agent-status.sh:9--17} 定义了只读状态查询接口。\verified\ 本报告源文件是本次任务的分析产物，PDF 只有在后文编译和版式检查完成后才视为有效交付。

\section{整体框架}
\label{sec:architecture}
configure 的 agent enhancement 由四层组成：入口层按调用名把 \code{cl/op/co} 分发到三条 CLI 路径；共享运行时层负责运行目录、锁、状态、事件和上下文；适配层负责每个 CLI 的 headless/resume 链；hook 层负责生命周期适配、路径保护和只读质量门禁。入口和适配的函数位置见
\code{bin/agent.sh:178--478}、\code{bin/agent.sh:480--834} 和
\code{bin/agent.sh:835--1229}。

\begin{figure}[htbp]
\centering
\begin{tikzpicture}[node distance=0.55cm and 0.7cm]
  \node[box] (input) {命令名与参数\\模型、输入、上限};
  \node[box,right=of input] (dispatch) {统一入口\\cl / op / co};
  \node[store,right=of dispatch] (runtime) {共享运行时\\run、lock、state、event};
  \node[box,below=of dispatch] (cli) {CLI 适配\\TUI 或 headless resume};
  \node[store,right=of cli] (data) {data/runs/<run-id>/\\可查询持久记录};
  \node[box,left=of cli] (hooks) {hooks\\preflight / guard / verify};
  \draw[arrow] (input) -- (dispatch);
  \draw[arrow] (dispatch) -- (runtime);
  \draw[arrow] (runtime) -- (cli);
  \draw[arrow] (cli) -- (data);
  \draw[relation] (hooks) -- (cli);
  \draw[relation] (hooks) |- (runtime);
\end{tikzpicture}
\caption{configure agent enhancement 的数据与控制关系。实线是主控制/持久化流，虚线是 hook 约束关系。}
\label{fig:configure-architecture}
\end{figure}

OMX 的层次更宽：CLI/skills 之上连接 session state、native hook、notifications 和 HUD；team 层再连接角色路由、任务 claim、worker mailbox、tmux pane、动态 scaling 与 worktree。其原生 hook 入口见
\code{codex-native-hook.ts:21888--22998}，team scaling 入口见
\code{scaling.ts:580--1921}；完整路径保留在参考源表中，因此两者不是同一规模的实现。

\section{关键细节}
\subsection{统一运行时状态机}
configure 先确定 workspace root，再创建数据根和运行目录；新运行写入 manifest/state 并加锁，resume 则先校验 schema、run id、agent、launcher 和 workspace，再读取历史 session/thread、模型、reasoning 和 variant。状态转移包括 created、running、finished、stopped、failed 和 blocked，回合开始、完成、失败及终态均追加 JSONL 事件。对应实现位于
\code{bin/agent-runtime.sh:238--360}、\code{bin/agent-runtime.sh:373--496} 和
\code{bin/agent-runtime.sh:550--615}。

锁使用 Bash 的 noclobber 创建；Unix 只在确认 owner PID 已退出时通过改名回收 stale lock，并在释放时再次核对 owner 是否为当前 PID。该设计处理了两个极端情形：活跃 owner 不可被误删，旧 owner 不会永久阻塞恢复。实现和回归证据分别为
\code{bin/agent-runtime.sh:126--167} 与 \code{bin/agent.test.sh:168--209}。

\subsection{上下文与 CLI 适配}
上下文发现沿当前目录到 Git 根逐层收集 \code{AGENTS.md}、\code{CODEX.md}、\code{CLAUDE.md} 和 \code{OPENCODE.md}，近处优先；运行时只把路径和行数写入 context 清单，正文仍由 agent 按需读取。实现见
\code{bin/agent-runtime.sh:292--360}，hook 的只读预检见 \code{hooks/codex-preflight.sh:36--71}。

Claude 通过 \code{claude -p} 和 \code{--resume} 绑定 session；OpenCode 通过持久日志提取 session id 并调用 \code{run -s}；Codex 通过 \code{codex exec --json} 提取 thread id 并调用 \code{exec resume}。三条链的关键入口分别为
\code{bin/agent.sh:345--463}、\code{bin/agent.sh:726--824} 和
\code{bin/agent.sh:1133--1209}。模型显式覆盖优先，resume 默认复用原记录中的模型及相关推理设置。

\subsection{hook 和质量边界}
\code{codex-hook.sh} 把 session-start、summary、before-edit、after-edit、stop 和 notify 统一映射；\code{--json} 通过 \code{codex-event.sh} 产生包含 schema、event、workspace、run/session/thread、turn 和 status 的 JSONL envelope。路径 guard 使用结构化 JSON 解析，拒绝仓库外、符号链接逃逸、日志和运行时数据路径。证据为
\code{hooks/codex-hook.sh:14--103}、\code{hooks/codex-event.sh:14--62} 和
\code{hooks/codex-guard.sh:97--225}。

\subsection{OMX 的关键扩展}
OMX 的 team 状态机由 \code{orchestrator.ts:8--146} 定义；角色提示和任务路由由
\code{role-router.ts:22--204} 提供；任务分配位于
\code{allocation-policy.ts:110--212}。worker runtime identity 的测试证明启动角色和模型 lane 会被保留，见
\code{worker-runtime-identity.test.ts}。\code{scaling.ts:109--276} 和
\code{scaling.ts:580--1430} 进一步将锁、pane receipt、进程 PID、Team owner、任务和 worktree 纳入 scale-up 事务；\code{worktree.ts:303--551} 管理 worktree 规划、创建、回滚和删除。

{\raggedright OMX 的 native hook 还处理 session owner、子 agent、终态取消、prompt provenance、状态重建和 replay dedupe；其关键入口可由
\code{codex-native-hook.ts:772--893}、\code{codex-native-hook.ts:1662--2188} 和
\code{codex-native-hook.ts:4005--4108} 看出。\code{process-identity.ts:17--79}、
\code{mode-binding-lease.ts:69--256} 和
\code{lifecycle-dedupe.ts:5--191} 分别覆盖进程身份、canonical session lease 和生命周期去重。\par}

\section{任务如何正确执行}
\label{sec:workflow}
本次正确执行路径如下。每一行都是一个可独立核验的动作，和实际脚本入口一一对应。

\begin{table}[htbp]
\centering
\caption{configure agent enhancement 的执行工作流}
\small
\begin{tabularx}{\textwidth}{@{}>{\raggedright\arraybackslash}p{0.06\textwidth}>{\raggedright\arraybackslash}X>{\raggedright\arraybackslash}p{0.29\textwidth}@{}}
\toprule
步骤 & 动作与停止条件 & 实现证据 \\
\midrule
1 & 解析 launcher、模型和驱动参数；非法数字或时长立即失败。 & \code{bin/agent.sh:250--280} \\
2 & 建立 workspace root、data 根和层级上下文清单。 & \code{bin/agent-runtime.sh:238--360} \\
3 & 新建或恢复 run；恢复前身份校验通过才取得锁。 & \code{bin/agent-runtime.sh:373--496} \\
4 & 绑定 CLI session/thread；写入 manifest/state 和 session-bound 事件。 & \code{bin/agent-runtime.sh:531--539} \\
5 & 运行首回合；每回合写 turn-start，完成或失败后更新计数。 & \code{bin/agent-runtime.sh:550--571} \\
6 & 检查 stop-file、max-turns、max-runtime 和连续失败阈值。 & \code{bin/agent-runtime.sh:574--596} \\
7 & 到达终态后追加 session-end/session-stop，并仅释放当前 owner 锁。 & \code{bin/agent-runtime.sh:598--615} \\
8 & 用 status、hook verify 和测试脚本检查持久输出及工作树。 & \code{bin/agent-status.sh:99--176}; \code{hooks/codex-verify.sh:157--289} \\
\bottomrule
\end{tabularx}
\end{table}

\inferred\ 这条顺序把“调用 CLI”和“管理运行”解耦：CLI 失败不会抹掉失败原因，外部程序也能在不执行 state 文件的情况下读取状态。对长期运行任务，这比单纯依赖工作目录下的临时日志更容易恢复和审计。

\section{实际执行情况}
\label{sec:execution}
configure 实际执行了 Shell 语法、agent fake CLI 回归、hook 回归、工具检查、coverage 报告检查、推荐插件安装检查、tag 链检查、Windows 静态契约检查和 Git 空白检查。configure 的核心回归命令为
\code{bash bin/agent.test.sh} 和 \code{bash hooks/hooks.test.sh}，均返回 0。

OMX 先执行 \code{npm ci} 和 \code{npm run build}，均返回 0；\code{npm ci} 报告 172 个包并提示 7 个 audit vulnerabilities（1 low、3 moderate、3 high）。随后运行三个聚焦套件：worker identity 与 cross-rebase 均通过；recent bug suite 在 \code{npm run build:runtime} 阶段因当前环境没有 \code{cargo} 返回 127。精确命令 \code{npm run test:explicit-terminal-contract} 没有定义，返回 1；实际存在的 compiled 命令运行后返回 1。

\section{实际输出情况}
\label{sec:outputs}
configure 的核心回归输出为 15 项 agent PASS 和 11 项 hook PASS。覆盖内容包括 Codex once/resume、manifest run id 错配、stale/busy/non-owner lock、跨 agent/workspace 恢复拒绝、OpenCode 文件输入、Claude session 绑定、max-turns、max-runtime、blocked、stop-file、JSON status、无工作目录垃圾和 hook 路径安全。

OMX worker identity 输出为 4 tests、4 pass、0 fail，耗时约 3.09 s；cross-rebase 输出为 3 tests、3 pass、0 fail，耗时约 2.87 s。compiled explicit terminal suite 的前置文件通过，但 native hook 文件最终失败；可复查的失败包括 malformed exec follow-up queue、persisted subagent resume guidance、session scope drift、unmatched sidecar blocker、native stop replay dedupe 和 permission-style auto-nudge。该结果不是 configure 的失败，表示 OMX 当前对照环境中的该套件未达到全绿。

\section{结果分析}
\label{sec:results}
\subsection{能力矩阵}
\begin{table}[htbp]
\centering
\caption{工作对象与 OMX 的 agent enhancement 对照}
\small
\begin{tabularx}{\textwidth}{@{}>{\raggedright\arraybackslash}p{0.23\textwidth}>{\raggedright\arraybackslash}X>{\raggedright\arraybackslash}X>{\raggedright\arraybackslash}p{0.17\textwidth}@{}}
\toprule
能力 & configure 最终版 & OMX HEAD & 证据等级 \\
\midrule
统一入口 & cl/op/co，共享 Shell runtime；Windows batch 转 PowerShell & 以 Codex 为执行核心的 npm CLI 与 plugin/setup 体系 & \verified\ / \code{bin/agent.sh:178--1237}; \code{package.json:1--70} \\
会话恢复 & run id + CLI session/thread，身份校验和参数复用 & canonical session、owner PID、子 agent/root 及更细的 session provenance & \verified\ 两边均有，OMX 更完整 \\
状态耐久性 & 原子替换 manifest/state，JSONL 事件，Unix stale lock & 原子状态、lease、reconciliation、derived state 和更严格的终态契约 & \verified\ / \inferred \\
运行边界 & max-turns、max-runtime、stop-file、blocked & explicit terminal contract、Stop hook 和 Autopilot/Ralph 状态机 & \verified\ 两边均有，语义不同 \\
上下文 & 从当前目录到仓库根分层发现四类说明文件 & scoped AGENTS、role/tier/posture、prompt composition & \verified\ / \inferred \\
hook & 轻量显式 adapter、JSON envelope、路径 guard/verify & native lifecycle hook、通知、dedupe、MCP、插件扩展 & \verified\ \\
多 worker & 未实现持久 mailbox、tmux authority、动态 scaling & team、任务 claim、worker mailbox、tmux pane、scale-up/down、worktree & \verified\ \\
跨平台 & Unix 与 Windows runtime 语义对齐；Windows 实机未验证 & macOS/Linux 推荐 tmux，Windows 有受限路径 & \verified\ / \unverified \\
依赖与部署 & Bash + 原生 CLI，低依赖，可直接放入配置仓库 & Node/npm，部分 runtime/explore 需要 Rust/Cargo & \verified\ \\
测试规模 & 15 agent + 11 hook 核心回归，静态 Windows 契约 & 443 个 test files；聚焦团队测试通过，native hook 聚焦套件未全绿 & \verified\ \\
\bottomrule
\end{tabularx}
\end{table}

\subsection{configure 的已证实优势}
第一，统一 launcher 让三种 CLI 共用运行边界和持久布局，减少了用户为不同工具记忆不同恢复方式的成本。第二，数据写入默认落在
\code{${HOME}/configure/data}，工作目录不再产生新版 \code{.agent.*} 日志；这对配置仓库和多 workspace 场景尤其重要。第三，恢复身份校验覆盖 run id、agent、launcher 和 workspace，且模型参数有“原记录复用、显式参数优先”的确定规则。第四，Unix stale lock 的改名回收和 owner 二次核验直接覆盖了并发/异常退出的高风险边界。上述判断均由
\code{data/README.md:7--27}、\code{bin/agent-runtime.sh:126--193} 和
\code{bin/agent.test.sh:120--325} 直接支持。

\subsection{工作对象的劣势及其根因}
主要劣势是规模和语义深度，而不是当前核心回归缺陷。\code{bin/agent.sh} 已达 1237 行，Windows runtime 达 1650 行，Unix 和 Windows 仍存在部分参数/驱动逻辑重复；这会增加后续修复需要同步两侧的风险。其次，状态是单 run 级别的轻量文件状态，没有 OMX 那样的 owner identity、pane authority、worker claim、mailbox 和 worktree 事务；多 worker 协作若直接在其上叠加，容易把“进程活着”误当成“任务仍归属该 worker”。第三，hook 是显式调用适配层，当前不能宣称已经接入某个 Codex 官方统一生命周期发现 API；\code{hooks/README.md:1--3} 已明确这个边界。第四，Windows 只有静态契约验证，独占句柄、PowerShell 解析器和真实 CLI 链尚未实机验证。

\section{综合分析}
\label{sec:synthesis}
从项目思路看，configure 选择的是“把成熟 CLI 当执行引擎，在配置仓库增加一个薄的、可移植的运行控制面”。共享 runtime 用稳定的文件协议连接三种 agent，hook 用显式命令降低对外部 API 版本的依赖，数据目录隔离则把可观测性从工作区副作用中分离出来。这解释了它在 HPC 登录节点、没有 Node/Rust 的机器和简单无人值守任务上的实际优势。

OMX 选择的是“把 Codex 扩展成完整的工作流和多 agent 平台”。其角色/tier/posture prompt composition、native agent catalog、terminal contract、session reconciliation、worker mailbox、tmux authority、动态 scaling、worktree 和插件/hook 生态互相咬合；这些能力带来更高的协作上限，也带来 Node、tmux、状态协议和更大的测试/维护成本。\code{README.md:36--42}、\code{README.md:121--167} 以及上述 team/state 源码共同支持这一判断。

两者的核心差异可以写成状态空间的差别：configure 主要维护
\begin{equation}
  S_{\rm cfg}=({\rm run},\ {\rm cli\ session},\ {\rm limits},\ {\rm lock},\ {\rm events}),
\end{equation}
而 OMX 还维护
\begin{equation}
  S_{\rm omx}=S_{\rm cfg}+({\rm owner},\ {\rm pane},\ {\rm worker},\ {\rm task\ claim},\ {\rm mailbox},\ {\rm worktree},\ {\rm hook\ provenance}).
\end{equation}
\inferred\ 第二个状态空间严格包含第一个的协作维度，因此 configure 在单进程/轻量部署上可以更简单、更快收敛，却不能据此声称在完整 team orchestration 上超过 OMX。

\section{结尾总结}
本次优化把 configure 从“按名称启动 CLI 的脚本集合”推进为“有统一运行协议的三 agent launcher”：运行创建、恢复、锁、上下文、回合上限、终态、事件和状态查询均有明确文件接口；Unix 与 Windows 有对应实现；hook 具备 JSON envelope 和 Git 路径质量边界。核心 configure 回归全部通过，OMX 的两个团队聚焦回归也通过。

最终判断是分场景的：在低依赖、跨 Codex/OpenCode/Claude、Unix/Windows 统一入口、显式停止和可观测性这组指标上，configure 已形成清晰优势；在多 worker 的身份权威、tmux 协作、动态伸缩、worktree 事务和 native hook 状态机上，OMX 仍然领先。报告不把未实现能力写成已完成能力。

\section{结尾评价}
\begin{table}[htbp]
\centering
\caption{优化评价与优先级}
\small
\begin{tabularx}{\textwidth}{@{}>{\raggedright\arraybackslash}p{0.22\textwidth}>{\raggedright\arraybackslash}X>{\raggedright\arraybackslash}p{0.22\textwidth}@{}}
\toprule
维度 & 当前评价 & 后续优先级 \\
\midrule
可靠性 & 恢复身份、锁 owner、运行上限和终态已有回归证据。 & 继续增加进程身份 token 和状态并发故障注入。 \\
可维护性 & 共享 runtime 已抽出；Unix/Windows 仍有较大重复。 & 把参数 schema 和事件字段抽为跨平台契约，减少双写。 \\
可观测性 & manifest/state/events/status 构成完整轻量闭环。 & 增加事件 schema 版本迁移和 redaction 检查。 \\
部署性 & Bash + 原生 CLI，适合没有 Node/Rust 的工作站/HPC。 & 在 Windows runner 和真实 CLI 上完成实机矩阵。 \\
协作上限 & 单 run 边界可靠，多 worker 编排能力不足。 & 若需求扩大，再引入 mailbox/worker claim/worktree，而不是把状态字段无限堆到单 run。 \\
\bottomrule
\end{tabularx}
\end{table}

\section{补充说明}
\label{sec:limits}
本报告有四类限制。其一，OMX 的 \code{test:recent-bug-regressions} 未能进入测试阶段，因为当前环境没有 \code{cargo}；不能据此判断 Rust runtime 本身通过或失败。其二，OMX 的 compiled explicit terminal suite 返回 1，报告只记录实际失败，不把未通过场景归因成单一原因。其三，Windows PowerShell parser、独占锁和真实 Codex/OpenCode/Claude CLI 没有在当前 macOS 上实机验证；Windows 结论限于静态契约。其四，已测的是无网络帮助入口启动时间，未测模型调用吞吐、首回合延迟或恢复延迟；因此启动数据只说明本机轻量入口开销，不能写成统一性能定律。

此外，configure 的 hook 不自动修改 Codex 配置或 \code{core.hooksPath}；启用方式仍由调用方 runner 决定，见
\pathref{hooks/README.md:39--51}。运行时 JSONL 的安全性依赖当前 Shell escape 实现，后续若允许任意非 ASCII 或更复杂的控制字符，应增加 jq/Python 解析回读测试。

\section{参考源}
\label{sec:sources}
\begin{longtable}{@{}>{\raggedright\arraybackslash}p{0.39\textwidth}>{\raggedright\arraybackslash}p{0.55\textwidth}@{}}
\toprule
来源 & 用途与关键行号 \\
\midrule
\endhead
\code{AGENTS.md:1--45} & 仓库边界、目录职责、测试惯例和 agent 入口。 \\
\code{bin/AGENTS.md:29--47} & launcher、数据目录、参数和跨平台约定。 \\
\code{bin/agent.sh:178--1237} & Claude/OpenCode/Codex 分发、驱动和参数解析。 \\
\code{bin/agent-runtime.sh:62--615} & manifest、state、lock、context、事件、恢复和终态。 \\
\code{bin/agent-status.sh:9--176} & 只读状态查询和 JSON 输出。 \\
\code{hooks/codex-hook.sh:14--103} & hook 事件映射和 JSON 模式。 \\
\code{hooks/codex-event.sh:14--62} & 统一 JSONL event envelope。 \\
\code{hooks/codex-guard.sh:97--225} & 仓库边界、符号链接和运行数据保护。 \\
\code{hooks/codex-preflight.sh:36--71} & 分层上下文只读预检。 \\
\code{hooks/codex-verify.sh:157--289} & Git、Shell、日志和数据路径验证。 \\
\code{bin/agent.test.sh:120--325} & configure agent fake CLI 回归。 \\
\code{hooks/hooks.test.sh:1--227} & hook 质量门禁回归。 \\
\code{data/README.md:7--27} & 运行时数据协议和隔离规则。 \\
\code{refer/git-rep/oh-my-codex/package.json:1--70} & OMX 版本、构建、测试和依赖入口。 \\
\code{refer/git-rep/oh-my-codex/src/team/orchestrator.ts:8--146} & team 状态机与 phase agent。 \\
\code{refer/git-rep/oh-my-codex/src/team/scaling.ts:109--1430} & claim lock、pane authority、scale-up 事务。 \\
\code{refer/git-rep/oh-my-codex/src/team/worktree.ts:303--551} & worktree 规划、创建、回滚和删除。 \\
\code{refer/git-rep/oh-my-codex/src/state/process-identity.ts:17--79} & 进程 owner token 和启动身份。 \\
\code{refer/git-rep/oh-my-codex/src/state/mode-binding-lease.ts:69--256} & canonical mode binding lease。 \\
\code{refer/git-rep/oh-my-codex/src/notifications/lifecycle-dedupe.ts:5--191} & lifecycle 通知去重。 \\
\code{refer/git-rep/oh-my-codex/src/scripts/codex-native-hook.ts:772--22998} & native hook 映射、session、终态和分发。 \\
\code{refer/git-rep/oh-my-codex/src/team/__tests__/worker-runtime-identity.test.ts} & worker runtime identity 聚焦测试。 \\
\code{refer/git-rep/oh-my-codex/src/team/__tests__/cross-rebase-smoke.test.ts} & cross-rebase 聚焦测试。 \\
\code{hooks/README.md:1--51} & hook API 边界和显式启用方式。 \\
\bottomrule
\end{longtable}

\section{附录与附件}
\label{sec:appendix}
\subsection{A：定量规模}
截至对照 HEAD，configure 全仓库 Git 跟踪文件数为 3271；本次新增/修改的 agent 相关文件共 17 个，按交付集合统计 5139 行，其中 \code{agent.sh} 1237 行、Unix runtime 616 行、PowerShell runtime 1650 行。OMX 的 \code{src/} 共 853 个文件、446080 行，其中 443 个测试文件；非测试源码文件 410 个。OMX team/state/notify/CLI 选定核心集合约 44863 行。此处行数是工程规模指标，不是质量分数。

\subsection{B：测试命令与退出码}
\small
\begin{longtable}{@{}>{\raggedright\arraybackslash}p{0.62\textwidth}>{\raggedright\arraybackslash}p{0.13\textwidth}>{\raggedright\arraybackslash}p{0.17\textwidth}@{}}
\toprule
命令 & 退出码 & 实际结果 \\
\midrule
\endhead
\code{bash -n bin/*.sh hooks/*.sh tools/*.sh plugins/*.sh skills/*/*.sh} & 0 & 通过 \\
\code{bash bin/agent.test.sh} & 0 & 15 PASS \\
\code{bash hooks/hooks.test.sh} & 0 & 11 PASS \\
\code{bash tools/configure-check.test.sh} & 0 & 通过；原有 plugins warning，无 error \\
\code{bash tools/coverage-report.test.sh} & 0 & 通过 \\
\code{bash plugins/install-recommended.test.sh} & 0 & 通过 \\
\code{bash skills/tag/tag-chain.test.sh} & 0 & 通过 \\
\code{bash bin/agent-windows.test.sh} & 0 & 静态契约通过；实机跳过 \\
\code{git diff --check} & 0 & 通过 \\
\code{npm ci}（OMX） & 0 & 172 packages；7 vulnerabilities \\
\code{npm run build}（OMX） & 0 & TypeScript 构建通过 \\
\code{npm run test:team:worker-runtime-identity}（OMX） & 0 & 4/4 通过 \\
\code{npm run test:team:cross-rebase-smoke}（OMX） & 0 & 3/3 通过 \\
\code{npm run test:recent-bug-regressions}（OMX） & 127 & 缺少 cargo，停在 build:runtime \\
\code{npm run test:explicit-terminal-contract}（OMX） & 1 & script 不存在 \\
\code{npm run test:explicit-terminal-contract:compiled}（OMX） & 1 & native hook 聚焦套件未全绿 \\
\bottomrule
\end{longtable}
\normalsize

无网络帮助入口启动的 3 次测量为：configure 使用
\code{bash bin/co --help}，耗时 \metric{0.08--0.09 s}，退出码均为 0；OMX 使用
\code{node dist/cli/omx.js --help}，耗时 \metric{0.29--0.33 s}，退出码均为 0。
该测量包含独立进程启动和参数帮助路径，不包含模型请求、网络、首回合、session 恢复或 team 建立；它支持“低依赖入口更轻”的局部判断，不能替代完整性能基准。

\subsection{C：交付前检查记录}
报告编译应使用以下闭环：在临时目录以 XeLaTeX 编译两遍，确认 PDF 非空；用可用的 PDF 工具检查页数和文本；若工具链不完整，明确写入未验证状态。\code{xelatex} 和 Ghostscript 在当前环境可用，\code{pdfinfo}、\code{pdftoppm}、\code{mutool} 和 \code{qpdf} 未发现，因此本次页数由 XeLaTeX/PDF 文件和文本提取检查确认，逐页图像渲染指标若无法取得则标为未验证，不冒充人工视觉通过。

\subsection{D：后续可验收项}
\begin{enumerate}\small
  \setlength{\itemsep}{0.1em}
  \item 在 Windows 或安装 PowerShell 的 CI 上运行 \code{bin/agent-windows.test.sh} 的实机分支，覆盖文件句柄锁、resume 和三类真实 CLI。
  \item 将 Unix/PowerShell 的参数和 event schema 生成自同一份契约，消除两侧重复解析。
  \item 若需要多 worker，再以独立的 worker identity、task claim、mailbox 和 worktree 事务为一组引入，并为每个状态转移增加故障注入测试。
  \item 在同一机器、同一 fake CLI 和相同输入下重复测量启动、首回合、恢复和状态查询延迟，再决定是否需要性能优化。
\end{enumerate}

\end{document}
